\section{Hermitization and the superlogarithmic bridge}\label{sec:hermitization}

We combine \cref{thm:intermediate,prop:row-replacement} with the smallest-singular-value inputs and the optimal sparse i.i.d. circular law. Throughout this section, fix an arbitrary deterministic integer sequence $d=d(n)$ satisfying the degree assumptions of the result under consideration, and set $p_n=d(n)/n$. Every assertion is made for this fixed sequence. The Bernoulli logarithmic-potential interface below is valid for every fixed $z\ne0$, so no sequence-dependent exceptional set of shifts is needed.

\subsection{A logarithmic-tail lemma}

\begin{lemma}[Negative logarithmic uniform integrability]\label{lem:negative-log-tail}
Let
\[
  \nu_n=\frac1n\sum_{j=1}^n\delta_{s_{j,n}}
\]
be a random probability measure on $[0,\infty)$, where $s_{1,n}\ge\cdots\ge s_{n,n}\ge0$. Suppose that, for a deterministic probability measure $\nu$,
\begin{enumerate}[label=\textup{(\roman*)}]
\item $\nu_n$ converges weakly to $\nu$ in probability;
\item
\[
  \int\log t\,\dd\nu_n(t)
  \xrightarrow{\Pp}
  \int\log t\,\dd\nu(t)\in\R;
\]
\item
\[
  \int t^2\,\dd\nu_n(t)=\bigOP.
\]
\end{enumerate}
Then, for every $\eps>0$,
\begin{equation}\label{eq:negative-log-UI}
  \lim_{T\to\infty}\limsup_{n\to\infty}
  \Pp\set{
    \int_{\set{\logm t>T}}\logm t\,\dd\nu_n(t)>\eps
  }=0.
\end{equation}
Consequently, if $r_n=o(n)$, then
\begin{equation}\label{eq:bottom-r-log-tail}
  \frac1n\sum_{j=n-r_n+1}^n\logm s_{j,n}
  \xrightarrow{\Pp}0.
\end{equation}
\end{lemma}

\begin{proof}
We first control positive logarithms. For $R\ge e$, put
\[
  c_R:=\sup_{t\ge R}\frac{\log t}{t^2},
\]
so $c_R\to0$. Then
\[
  \int_{t>R}\log t\,\dd\nu_n(t)
  \le c_R\int t^2\,\dd\nu_n(t),
\]
which tends to zero in the iterated limit $n\to\infty$, $R\to\infty$, in probability. Weak convergence applied to $\min(\logp t,\log R)$ therefore gives
\begin{equation}\label{eq:positive-log-convergence}
  \int\logp t\,\dd\nu_n(t)
  \xrightarrow{\Pp}
  \int\logp t\,\dd\nu(t).
\end{equation}

The bounded-in-probability second moments also imply
\[
  \int t^2\,\dd\nu(t)<\infty.
\]
Indeed, for each $L<\infty$, weak convergence in probability gives
\[
  \int(t^2\wedge L)\,\dd\nu_n(t)
  \xrightarrow{\Pp}
  \int(t^2\wedge L)\,\dd\nu(t).
\]
Choose $M<\infty$ such that
\[
  \limsup_{n\to\infty}
  \Pp\set{\int t^2\,\dd\nu_n(t)>M}<\frac12.
\]
Because the truncated integral is bounded above by the full second moment, its deterministic probability limit cannot exceed $M$. Hence
\[
  \int(t^2\wedge L)\,\dd\nu(t)\le M
\]
for every $L$, and monotone convergence proves the claim.

Since $\log t=\logp t-\logm t$, assumption~(ii) and \eqref{eq:positive-log-convergence} imply
\begin{equation}\label{eq:negative-log-total}
  \int\logm t\,\dd\nu_n(t)
  \xrightarrow{\Pp}
  \int\logm t\,\dd\nu(t)<\infty.
\end{equation}
For fixed $T$, the function
\[
  f_T(t)=\min(\logm t,T),
  \qquad f_T(0)=T,
\]
is bounded and continuous. Thus
\[
  \int f_T\,\dd\nu_n
  \xrightarrow{\Pp}
  \int f_T\,\dd\nu.
\]
Subtracting from \eqref{eq:negative-log-total} yields
\[
  \int(\logm-f_T)\,\dd\nu_n
  \xrightarrow{\Pp}
  \int(\logm-f_T)\,\dd\nu.
\]
The deterministic right-hand side tends to zero by monotone convergence. On $\set{\logm t>2T}$,
\[
  \logm t\le2(\logm t-T)=2(\logm t-f_T(t)).
\]
This proves \eqref{eq:negative-log-UI}, after renaming $2T$ as $T$.

Finally, for fixed $T$,
\[
  \frac1n\sum_{j=n-r_n+1}^n\logm s_{j,n}
  \le T\frac{r_n}{n}
  +\int_{\set{\logm t>T}}\logm t\,\dd\nu_n(t).
\]
Let $n\to\infty$ and then $T\to\infty$.
\end{proof}

\subsection{The Bernoulli hermitization interface}

Let $B_n$ denote an $n\times n$ matrix with i.i.d. $\Ber(d/n)$ entries. For fixed $z\in\C$, write
\[
  H_{n,z}:=\frac{B_n}{\sqrt d}-zI_n,
  \qquad
  \nu^B_{n,z}:=\frac1n\sum_{j=1}^n\delta_{s_j(H_{n,z})}.
\]
Let $\nu_z^G$ denote the deterministic limiting singular-value law of
$n^{-1/2}G_n-zI_n$, where $G_n$ has i.i.d. standard complex Gaussian entries, and put
\begin{equation}\label{eq:Gaussian-log-potential}
  U^\circ(z):=-\int_\C\log\abs{z-w}\,\dd\mu_{\mathrm{circ}}(w)
  =-\int_0^\infty\log t\,\dd\nu_z^G(t).
\end{equation}

\begin{proposition}[Bernoulli hermitization interface]\label{prop:Bernoulli-hermitization}
Assume $d\to\infty$ and $d=o(n)$. For every fixed $z\in\C\setminus\set{0}$:
\begin{enumerate}[label=\textup{(\roman*)}]
\item $\nu^B_{n,z}\Rightarrow\nu_z^G$ in probability;
\item
\begin{equation}\label{eq:Bernoulli-full-log}
  -\int_0^\infty\log t\,\dd\nu^B_{n,z}(t)
  \xrightarrow{\Pp}U^\circ(z),
\end{equation}
equivalently, the full untruncated logarithmic integral converges to the finite value
$\int\log t\,\dd\nu_z^G(t)$;
\item
\begin{equation}\label{eq:Bernoulli-second-moment}
  \int_0^\infty t^2\,\dd\nu^B_{n,z}(t)=\bigOP.
\end{equation}
\end{enumerate}
\end{proposition}

\begin{proof}
Part~(i) is \cite[Fact~10.2]{SSS25}, applied to the deterministic sequence $p_n=d/n$ and to the Bernoulli model. The Gaussian limit in that fact is $\nu_z^G$. By \cite[Fact~10.1]{SSS25}, the measure $\nu_z^G$ is compactly supported and satisfies
\[
  \nu_z^G([0,t))\le Ct,
  \qquad t>0,
\]
and \eqref{eq:Gaussian-log-potential} holds. The mass bound at zero and compact support imply that the limiting logarithmic integral is finite.

We reproduce the squeeze that gives the full, untruncated convergence in part~(ii). Define
\begin{equation}\label{eq:Bernoulli-Un}
  U_n(z):=-\frac1n\sum_{j=1}^n\log s_j(H_{n,z}),
\end{equation}
with the extended value $+\infty$ when $H_{n,z}$ is singular. Choose a deterministic sequence $\eps_n\downarrow0$ sufficiently slowly for the estimates in Sections~8--10 of \cite{SSS25}, and write $\eps=\eps_n$. Set
\[
  m=\floor{(1-\eps)n},
  \qquad
  \ell=\floor{(1-\eps/4)m},
  \qquad
  a=\ceil{2\eps n},
  \qquad
  b=\floor{(1-\eps)n}.
\]
Let $H_{m,z}$ be the top-left $m\times m$ principal minor of $H_{n,z}$. Thus its entries retain the normalization $d^{-1/2}$ used for $H_{n,z}$. Define
\begin{align}
  T_n(z)&:=-\frac1n\sum_{j=1}^{\ell}\log s_j(H_{m,z}),
  \label{eq:SSS-T}\\
  T_n^{(1)}(z)&:=-\frac1n\sum_{j=1}^{\ell}\log s_j(H_{n,z}),
  \label{eq:SSS-T1}\\
  T_n^{(2)}(z)&:=-\frac1n\sum_{j=a}^{b}\log s_j(H_{n,z})
  -\frac{\ell-\floor{(1-3\eps)n}}{n}\log s_{\ell}(H_{m,z}).
  \label{eq:SSS-T2}
\end{align}
These are the integer-rounded versions of the quantities in Sections~8--10 of \cite{SSS25}.

For every fixed $z\ne0$, \cite[Lemma~8.2]{SSS25} gives
\begin{equation}\label{eq:SSS-upper-squeeze}
  U_n(z)\le T_n(z)+\littleoP.
\end{equation}
The first deterministic comparison in \cite[Fact~9.1]{SSS25} gives
\begin{equation}\label{eq:SSS-T-upper}
  T_n(z)\le T_n^{(2)}(z).
\end{equation}
For completeness, we derive the lower comparison directly from monotonicity of the ordered singular values. Put
\[
  x_j=-\log s_j(H_{n,z}),\qquad 1\le j\le n.
\]
Since the singular values are nonincreasing, $x_1\le\cdots\le x_n$. Hence
\begin{equation}\label{eq:SSS-U-lower-exact}
  U_n(z)=\frac1n\sum_{j=1}^n x_j
  \ge\frac1\ell\sum_{j=1}^{\ell}x_j
  =\frac n\ell T_n^{(1)}(z).
\end{equation}
Moreover,
\[
  \frac\ell n=(1-\eps/4)(1-\eps)+O(n^{-1}).
\]
Finally, \cite[Lemma~9.3]{SSS25}, proved there from Facts~10.1--10.2 and Lemmas~10.3--10.5, yields
\begin{equation}\label{eq:SSS-truncated-convergence}
  T_n^{(1)}(z)\xrightarrow{\Pp}U^\circ(z),
  \qquad
  T_n^{(2)}(z)\xrightarrow{\Pp}U^\circ(z)
\end{equation}
when $\eps_n\downarrow0$ sufficiently slowly. Since $n/\ell\to1$, \eqref{eq:SSS-U-lower-exact} and the first convergence in \eqref{eq:SSS-truncated-convergence} give
\[
  U_n(z)\ge U^\circ(z)-\littleoP.
\]
On the other hand, \eqref{eq:SSS-upper-squeeze}, \eqref{eq:SSS-T-upper}, and the second convergence in \eqref{eq:SSS-truncated-convergence} give
\[
  U_n(z)\le U^\circ(z)+\littleoP.
\]
Therefore $U_n(z)\to U^\circ(z)$ in probability. In particular,
\[
  \Pp\set{H_{n,z}\text{ is singular}}\to0.
\]
Redefining $U_n(z)$ arbitrarily on that event does not change the conclusion. This proves \eqref{eq:Bernoulli-full-log} for every fixed $z\ne0$.

For part~(iii),
\begin{align*}
  \int t^2\,\dd\nu^B_{n,z}(t)
  &=\frac1n\norm{H_{n,z}}_{\HS}^2\\
  &\le \frac{2}{nd}\sum_{i,j=1}^n(B_n)_{ij}+2\abs z^2.
\end{align*}
The sum has mean $nd$ and is $O_{\Pp}(nd)$, proving \eqref{eq:Bernoulli-second-moment}.
\end{proof}

\begin{corollary}[Bernoulli bottom tail]\label{cor:Bernoulli-bottom-tail}
Fix $z\in\C\setminus\set{0}$. For every deterministic integer sequence $r_n=o(n)$,
\begin{equation}\label{eq:Bernoulli-bottom-tail}
  \frac1n\sum_{j=n-r_n+1}^n
  \logm s_j\!\left(\frac{B_n}{\sqrt d}-zI_n\right)
  =\littleoP.
\end{equation}
\end{corollary}

\begin{proof}
Apply \cref{lem:negative-log-tail} to the three conclusions of
\cref{prop:Bernoulli-hermitization}.
\end{proof}

\subsection{Lower-band regularized comparison of a large prefix}

Fix $\sigma\in(0,1)$, assume \eqref{eq:lower-band}, fix $z\in\C\setminus\set{0}$, and define the unnormalized shifted matrices
\begin{equation}\label{eq:shifted-comparators}
  A=M_n-z\sqrt d\,I_n,
  \qquad
  \wt A=B_n-z\sqrt d\,I_n.
\end{equation}
Here $B_n$ is independent of $M_n$ and has i.i.d. $\Ber(p)$ entries, $p=d/n$.

By \cref{prop:LLY-smallest}, for a sufficiently large absolute constant $C_L$,
\begin{equation}\label{eq:lower-tau}
  \tau_n=n^{-C_L\log\log n}
\end{equation}
satisfies
\begin{equation}\label{eq:lower-good-event}
  \Pp(\mathcal L_n(z)^c)=O(d^{-1/2})=o(1),
  \qquad
  \mathcal L_n(z)=\set{s_n(A)\ge\tau_n}\cap\set{s_n(\wt A)\ge\tau_n}.
\end{equation}
The shift restriction in \cref{prop:LLY-smallest} holds because
$\abs z\sqrt d\le\sqrt d\log\log d$ for fixed $z\ne0$ and all sufficiently large $n$.

Choose
\begin{equation}\label{eq:lower-prefix-parameters}
  r=\ceil{\frac n{\sqrt d}},
  \qquad
  k=n-r,
  \qquad
  \eta=\frac{\tau_n}{n}.
\end{equation}
Let $C_M$ and $C_B$ be the first $k$ rows of $A$ and $\wt A$, respectively. In \cref{prop:row-replacement},
\[
  Q_n=p(1-p)(r+1)\asymp\sqrt d,
\]
and the deterministic shifts are $u_i=-z\sqrt d\,e_i$. Thus
\[
  U\le\sqrt{\eta^2+(\sqrt n+\abs z\sqrt d)^2}\le2(1+\abs z)\sqrt n
\]
and
\begin{equation}\label{eq:lower-U-over-eta}
  \log\frac U\eta=O(\log n\,\log\log n)=O(d\log d).
\end{equation}
Consequently,
\begin{align}
  \frac1n\E\abs{\Psi_\eta(C_M)-\Psi_\eta(C_B)}
  &\le O\!\left(d^{-1/4}\sqrt{\log d}\right)
  +O\!\left(d^{3/2}\log d\,e^{-c\sqrt d}\right)
  \notag\\
  &=o(1).
  \label{eq:lower-regularized-comparison}
\end{align}
Markov's inequality gives
\begin{equation}\label{eq:lower-regularized-op}
  \frac1n\abs{\Psi_\eta(C_M)-\Psi_\eta(C_B)}=\littleoP.
\end{equation}

\subsection{Removing the lower-band regularization}

Let $t_1(C)\ge\cdots\ge t_k(C)$ be the singular values of a $k\times n$ prefix. Cauchy interlacing gives
\[
  t_i(C_M)\ge s_{i+r}(A)\ge s_n(A),
\]
and similarly for $C_B$. On $\mathcal L_n(z)$, every prefix singular value is at least $\tau_n$. Hence
\begin{align*}
  0\le\Psi_\eta(C)-\log\vol(C)
  &=\frac12\sum_{i=1}^k\log\left(1+\frac{\eta^2}{t_i(C)^2}\right)\\
  &\le\frac{k\eta^2}{2\tau_n^2}
  \le\frac1{2n}.
\end{align*}
After division by $n$, the error is at most $1/(2n^2)$. Combining this with
\eqref{eq:lower-regularized-op} and \eqref{eq:lower-good-event}, and extending both logarithmic-volume random variables by zero on $\mathcal L_n(z)^c$, gives
\begin{equation}\label{eq:lower-prefix-volume-comparison}
  \frac1n\abs{\log\vol(C_M)-\log\vol(C_B)}=\littleoP.
\end{equation}

\subsection{From lower-band prefix volumes to square determinants}

Let $H$ be an invertible $n\times n$ matrix, and let $C$ be its first $n-r$ rows. Interlacing implies
\[
  \sum_{j=r+1}^n\log s_j(H)
  \le\log\vol(C)
  \le\sum_{j=1}^{n-r}\log s_j(H).
\]
Therefore
\begin{equation}\label{eq:prefix-square-error}
  \abs{\log\abs{\det H}-\log\vol(C)}
  \le\sum_{j=1}^r\logp s_j(H)
  +\sum_{j=n-r+1}^n\logm s_j(H).
\end{equation}
For the positive logarithms, concavity gives
\begin{equation}\label{eq:positive-prefix-tail}
  \frac1n\sum_{j=1}^r\logp s_j(H)
  \le\frac r{2n}\log\left(1+\frac{\norm H_{\HS}^2}{r}\right).
\end{equation}
Since $\norm A_{\HS}^2\le2(1+\abs z^2)nd$ deterministically and
$\norm{\wt A}_{\HS}^2=O_{\Pp}(nd)$,
\begin{equation}\label{eq:lower-positive-tail}
  \frac1n\sum_{j=1}^r\logp s_j(A)
  =O\!\left(\frac{\log d}{\sqrt d}\right)=o(1),
\end{equation}
with the analogous $o_{\Pp}(1)$ estimate for $\wt A$.

For the negative row-slice tail, use bottom-up indexing
$a_j=s_{n-j+1}(A)$. On $\mathcal L_n(z)$ and the event in \cref{thm:intermediate},
\begin{equation}\label{eq:lower-bottom-q0}
  \frac1n\sum_{j=1}^{q_0}\logm a_j
  \le\frac{q_0}{n}\log\frac1{\tau_n}
  =O\!\left(\frac{\log n\,\log\log n}{d^2}\right)=o(1),
\end{equation}
and
\begin{equation}\label{eq:lower-bottom-intermediate}
  \frac1n\sum_{j=q_0+1}^{r}\logm a_j
  \le\frac rn C_{\isv}(z)d^{1/96}
  =O(d^{-47/96})=o(1).
\end{equation}
Thus
\begin{equation}\label{eq:lower-row-bottom-tail}
  \frac1n\sum_{j=n-r+1}^n\logm s_j(A)=\littleoP.
\end{equation}
By \cref{cor:Bernoulli-bottom-tail}, and because multiplication by $\sqrt d\ge1$ can only decrease $\logm$,
\begin{equation}\label{eq:lower-Bernoulli-bottom-tail}
  \frac1n\sum_{j=n-r+1}^n\logm s_j(\wt A)=\littleoP.
\end{equation}
Applying \eqref{eq:prefix-square-error} to $(A,C_M)$ and $(\wt A,C_B)$, and using \eqref{eq:lower-prefix-volume-comparison}, yields
\begin{equation}\label{eq:lower-unnormalized-det-comparison}
  \frac1n\log\abs{\det A}
  -\frac1n\log\abs{\det\wt A}
  \xrightarrow{\Pp}0.
\end{equation}
On the event where a determinant vanishes, define its normalized logarithm arbitrarily; by \eqref{eq:lower-good-event}, this changes a random variable only on an event of probability $o(1)$. Since the common factor $d^{-n/2}$ cancels,
\begin{equation}\label{eq:lower-normalized-det-comparison}
  \frac1n\log\abs{\det\!\left(\frac{M_n}{\sqrt d}-zI_n\right)}
  -\frac1n\log\abs{\det\!\left(\frac{B_n}{\sqrt d}-zI_n\right)}
  \xrightarrow{\Pp}0
\end{equation}
for every fixed $z\ne0$, hence for Lebesgue-almost every fixed shift.

\subsection{The lower-polylogarithmic circular law}

\begin{theorem}[Lower-polylogarithmic regime]\label{thm:lower-band-circular}
Fix $\sigma\in(0,1)$. There exists an absolute constant $C_*>0$ such that, whenever
\[
  C_*\log n\le d\le(\log n)^{1+\sigma},
\]
one has
\[
  \mu_{M_n/\sqrt{d(1-d/n)}}\xrightarrow{\Pp}\mu_{\mathrm{circ}}.
\]
The constant $C_*$ is independent of $\sigma$; the threshold from which the conclusion holds may depend on $\sigma$.
\end{theorem}

\begin{proof}
Set
\[
  X_n=\frac{M_n}{\sqrt d},
  \qquad
  Y_n=\frac{B_n}{\sqrt d}.
\]
Then
\[
  \frac1n\norm{X_n}_{\HS}^2=1
\]
deterministically and
\[
  \E\frac1n\norm{Y_n}_{\HS}^2=1.
\]
Thus the Hilbert--Schmidt condition in the replacement principle
\cite[Theorem~2.1]{TV10} is satisfied. Equation~\eqref{eq:lower-normalized-det-comparison} gives its logarithmic-determinant condition for every fixed $z\ne0$, and therefore for Lebesgue-almost every fixed $z$. No estimate uniform in $z$ and no union bound over shifts is required. Consequently,
\begin{equation}\label{eq:BL-row-Bernoulli}
  d_{\BL}(\mu_{X_n},\mu_{Y_n})\xrightarrow{\Pp}0.
\end{equation}
The optimal sparse i.i.d. circular law \cite{RT19,SSS25} applies because
\[
  p=d/n\to0,
  \qquad
  np=d\to\infty,
\]
and gives
\[
  d_{\BL}(\mu_{Y_n},\mu_{\mathrm{circ}})\xrightarrow{\Pp}0.
\]
The triangle inequality therefore yields
\begin{equation}\label{eq:sqrt-d-circular}
  \mu_{M_n/\sqrt d}\xrightarrow{\Pp}\mu_{\mathrm{circ}}.
\end{equation}

For every square matrix $H$, Schur triangularization gives
\[
  \sum_{j=1}^n\abs{\lambda_j(H)}^2\le\norm H_{\HS}^2.
\]
Consequently,
\[
  \int\abs z\,\dd\mu_H(z)
  \le\left(\frac1n\norm H_{\HS}^2\right)^{1/2}.
\]
Let $c_n=(1-d/n)^{-1/2}\to1$ and write $\mu_n=\mu_{M_n/\sqrt d}$. For any bounded $1$-Lipschitz function $f$,
\begin{align*}
  \abs{\int f(c_nz)\,\dd\mu_n(z)-\int f(z)\,\dd\mu_n(z)}
  &\le\abs{c_n-1}\int\abs z\,\dd\mu_n(z)\\
  &\le\abs{c_n-1}
  \left(\frac1n\norm{M_n/\sqrt d}_{\HS}^2\right)^{1/2}\\
  &=\abs{c_n-1}.
\end{align*}
Thus the bounded-Lipschitz distance between the ESDs at the two normalizations tends to zero deterministically. Combining this with \eqref{eq:sqrt-d-circular} proves the theorem.
\end{proof}

\subsection{The superlogarithmic bridge}\label{subsec:superlog-bridge}

\begin{proposition}[Superlogarithmic bridge]\label{prop:superlog-bridge}
Fix $\eps>0$. Let $d=d(n)$ be an integer sequence satisfying
\[
  (\log n)^{1+\eps}\le d,
  \qquad
  d=o(n).
\]
Then
\[
  \mu_{M_n/\sqrt{d(1-d/n)}}\xrightarrow{\Pp}\mu_{\mathrm{circ}}.
\]
\end{proposition}

\begin{proof}
For all sufficiently large $n$, one has $d\le n/2$. Let $B_\eps$ be supplied by \cref{cor:uniform-polynomial}, fix $z\in\C\setminus\set{0}$, and define $A,\wt A$ by \eqref{eq:shifted-comparators}. Put
\[
  \tau_{n,\eps}=n^{-B_\eps},
  \qquad
  \mathcal L_{n,\eps}(z)
  =\set{s_n(A)\ge\tau_{n,\eps}}
  \cap\set{s_n(\wt A)\ge\tau_{n,\eps}}.
\]
By \cref{cor:uniform-polynomial},
\begin{equation}\label{eq:superlog-good-event}
  \Pp(\mathcal L_{n,\eps}(z)^c)=O(d^{-1/2})=o(1).
\end{equation}
Choose the balanced exponent
\begin{equation}\label{eq:balanced-a}
  a=a_\eps:=\frac{2+\eps}{2(1+\eps)}.
\end{equation}
Then
\begin{equation}\label{eq:a-identities}
  1-a=\frac{\eps}{2(1+\eps)},
  \qquad
  a(1+\eps)-1=\frac\eps2.
\end{equation}
Set
\begin{equation}\label{eq:superlog-prefix-parameters}
  r=\ceil{\frac n{d^a}},
  \qquad
  k=n-r,
  \qquad
  \eta=\frac{\tau_{n,\eps}}n.
\end{equation}
Let $C_M^{(\eps)}$ and $C_B^{(\eps)}$ be the first $k$ rows of $A$ and $\wt A$. Since $p=d/n=o(1)$,
\begin{equation}\label{eq:superlog-Q}
  Q=p(1-p)(r+1)\asymp d^{1-a}
  \ge c(\log n)^{\eps/2}\to\infty.
\end{equation}
Moreover, $U\le2(1+\abs z)\sqrt n$ and
\[
  \log\frac U\eta=O_\eps(\log n).
\]
Therefore \cref{prop:row-replacement} gives
\begin{align}
  \frac1n\E\abs{
    \Psi_\eta(C_M^{(\eps)})-\Psi_\eta(C_B^{(\eps)})
  }
  &\le C\sqrt{\frac{\log(ed)}{d^{1-a}}}
  +C_\eps\sqrt d\,\log n\,e^{-cd^{1-a}}
  \notag\\
  &=o(1).
  \label{eq:superlog-replacement}
\end{align}
Indeed, $\log d=o(d^{1-a})$. Also
\[
  \log(\sqrt d\log n)=\tfrac12\log d+\log\log n
  =O_\eps(\log d)=o(d^{1-a}),
\]
where $d\ge(\log n)^{1+\eps}$ is used in the middle estimate. Thus the exponential term dominates its prefactor uniformly in the stated range.

Markov's inequality and the same interlacing calculation used in the lower band yield, after extending both log-volume random variables by zero on $\mathcal L_{n,\eps}(z)^c$,
\begin{equation}\label{eq:superlog-prefix-volume}
  \frac1n\abs{
    \log\vol(C_M^{(\eps)})-\log\vol(C_B^{(\eps)})
  }=\littleoP.
\end{equation}
Here the regularization is removed on $\mathcal L_{n,\eps}(z)$ exactly as before because $\eta=\tau_{n,\eps}/n$.

It remains to pass from prefixes to square determinants. By \eqref{eq:positive-prefix-tail},
\begin{equation}\label{eq:superlog-positive-tail}
  \frac1n\sum_{j=1}^r\logp s_j(A)
  =O(d^{-a}\log d)=o(1),
\end{equation}
with the analogous $o_{\Pp}(1)$ estimate for $\wt A$. On $\mathcal L_{n,\eps}(z)$, every one of the bottom $r$ row-slice singular values may pay the polynomial smallest-singular-value cost:
\begin{align}
  \frac1n\sum_{j=n-r+1}^n\logm s_j(A)
  &\le\frac rn B_\eps\log n
  \notag\\
  &\le C_\eps d^{-a}\log n
  \le C_\eps(\log n)^{-\eps/2}=o(1),
  \label{eq:superlog-row-bottom-tail}
\end{align}
where \eqref{eq:a-identities} is used in the last step. Since $r=o(n)$, \cref{cor:Bernoulli-bottom-tail} gives the corresponding estimate for $B_n/\sqrt d-zI_n$; multiplication by $\sqrt d\ge1$ can only decrease $\logm$, so the same estimate holds for $\wt A$.

Applying \eqref{eq:prefix-square-error} to both matrices and using \eqref{eq:superlog-prefix-volume}, we obtain
\begin{equation}\label{eq:superlog-det-comparison}
  \frac1n\log\abs{\det\!\left(\frac{M_n}{\sqrt d}-zI_n\right)}
  -\frac1n\log\abs{\det\!\left(\frac{B_n}{\sqrt d}-zI_n\right)}
  \xrightarrow{\Pp}0
\end{equation}
for every fixed $z\ne0$.

The Hilbert--Schmidt estimates, the replacement principle, and the optimal sparse i.i.d. circular law now give
\[
  \mu_{M_n/\sqrt d}\xrightarrow{\Pp}\mu_{\mathrm{circ}}
\]
exactly as in the proof of \cref{thm:lower-band-circular}. Finally, $d=o(n)$ implies $(1-d/n)^{-1/2}\to1$, and the same Schur/Hilbert--Schmidt argument transfers the conclusion to the variance normalization.
\end{proof}

\begin{remark}[Relation to the Li--Litvak--Yu range]
For every fixed $\eta\in(0,1)$, \cite[Theorem~1.3]{LLY26a} already gives the circular law when
\[
  d\ge(\log n)^{2+\eta},
  \qquad
  d=o(n).
\]
Thus, to connect the lower band to the previously available result, \cref{prop:superlog-bridge} is needed only on the intervening range between $(\log n)^{1+\eps}$ and a threshold of the form $(\log n)^{2+\eta}$. We stated and proved the bridge for the entire superlogarithmic range because the same argument is uniform there and requires no additional input.
\end{remark}

\subsection{Completion of the main theorem}

\begin{proof}[Proof of \cref{thm:main}]
Fix $\sigma=1/2$ and split the indices into two sets:
\[
  d\le(\log n)^{3/2}
  \qquad\text{and}\qquad
  d>(\log n)^{3/2}.
\]
On the first set, apply \cref{thm:lower-band-circular}. On the second set, apply \cref{prop:superlog-bridge} with $\eps=1/2$. Each subsequence converges in probability to the same deterministic circular law. Therefore the original sequence does as well.
\end{proof}
